% !PRINTED 439
\weberpage{439}
%<<CONT
zusammenhängen (\S.\,37, 9.). Es wird also durch $A$ nicht nur aus jedem Punkte ein Punkt, sondern auch aus jeder geraden Linie eine gerade Linie abgeleitet.

Die durch (1) ausgedrückte Beziehung der Punkte $y$ zu den Punkten $x$ kann als eine \gesperrt{Abbildung der Ebene in sich selbst} bezeichnet werden, wobei jedem Punkte ein bestimmter anderer Punkt und jeder geraden Linie eine gerade Linie entspricht. Eine solche Abbildung heisst in der Geometrie auch eine Collineation. Der Punkt $y$ heisst das \gesperrt{Bild} des Punktes $x$ und die gerade Linie $v$ das \gesperrt{Bild} der geraden Linie $u$. Hat man die Bilder zweier Punkte, so ist die Verbindungslinie dieser Bilder das Bild der Verbindungslinie der beiden Originalpunkte.

Bedeutet $S$ eine zweite ternäre lineare Substitution und
\[ A' = S^{-1}AS \]
die aus $A$ durch $S$ transformirte Substitution, so ist
\[ (y') = A'(x') \]
gleichbedeutend mit
\[ (y) = A(x), \]
wenn
\[ (y) = S(y'), \quad (x) = S(x') \]
gesetzt wird. Statt nun durch $S$ eine Abbildung zu definiren, kann man auch eine \gesperrt{Coordinatentransformation} darunter verstehen, indem man unter $x_1', x_2', x_3'$ die Coordinaten des Punktes $(x)$ in einem neuen, durch $S$ bestimmten Coordinatensysteme versteht. Dann wird der Zusammenhang zwischen den beiden Punkten $x$, $y$ ebenso wie durch $A$ auch durch die transformirte Substitution $A'$ ausgedrückt, und die Transformation der Substitutionen ist also gleichbedeutend mit der Transformation des Coordinatensystems. Die Composition der transformirten Substitutionen geschieht, wie wir schon früher gesehen haben, ebenso wie die der ursprünglichen, und es entsteht also durch Transformation aus jeder Gruppe eine isomorphe Gruppe.

Von grosser Wichtigkeit ist nun die Untersuchung der Punkte, die bei der Abbildung durch eine Substitution $A$ ungeändert bleiben.

Nehmen wir $A$ in der Form an:
\[ A = \begin{pmatrix} \alpha, & \beta, & \gamma \\ \alpha', & \beta', & \gamma' \\ \alpha'', & \beta'', & \gamma'' \end{pmatrix}, \]
%CONT>>
